% =============================================================
% ANEXO — Gestión del proyecto (GEP)
% Adaptado a partir de la Entrega 4 GEP (Fita Inicial) con las
% reconciliaciones necesarias respecto a la ejecución real:
%   - 28 → 23 emociones (filtrado documentado en cap 3 §3.1.1)
%   - 364 → 276 probes activos (cap 5 §5.1)
%   - G6 incorpora logit lens como 5ª técnica (cap 5 §5.2)
%   - T26 fine-tuning post-SVD: ejecutado, no opcional (cap 6 §6.4)
%   - Sección A.8 nueva: ejecución real vs planificación
% =============================================================

\section{Gestión del proyecto (GEP)}
\label{sec:gep}

Este anexo recoge la documentación de gestión del proyecto exigida por
el módulo GEP de la FIB. Se redactó originalmente como Entrega 4
(\textit{Fita Inicial}) durante el primer mes del TFG y se actualiza
aquí con las reconciliaciones necesarias respecto a la ejecución real
del proyecto. La sección final (\S\ref{sec:gep-ejecucion-real}) compara
la planificación inicial con lo efectivamente ejecutado.

% -------------------------------------------------------------
\subsection{Contextualización}
\label{sec:gep-contexto}
% -------------------------------------------------------------

La arquitectura Transformer~\cite{vaswani2017attention} ha redefinido el
procesamiento del lenguaje natural. BERT~\cite{devlin2019bert}, con
110~millones de parámetros, fue uno de los primeros modelos en demostrar
que una misma red pre-entrenada podía alcanzar resultados de referencia
en tareas tan distintas como clasificación de texto, extracción de
información o respuesta a preguntas. El problema es que 110~millones de
parámetros es hoy la talla \emph{pequeña}: modelos más recientes escalan
a cientos de miles de millones, y desplegarlos en entornos con recursos
limitados ---dispositivos móviles, sistemas embebidos, infraestructuras
con restricciones de latencia o energía--- sigue siendo un reto
abierto~\cite{strubell2019energy}.

Este TFG se centra en una técnica concreta de compresión ---la
descomposición en valores singulares (SVD)--- aplicada de forma
selectiva a las capas lineales de un modelo Transformer, y la combina
con cinco técnicas de interpretabilidad mecánica para entender
\emph{por qué} ciertos componentes toleran la compresión y otros no.

\subsubsection{Marco del proyecto en la FIB}

El proyecto se enmarca en el Grado en Ingeniería Informática de la
Facultat d'Informàtica de Barcelona (FIB-UPC), especialidad de
\textbf{Computación}, e integra contenidos de varias asignaturas:

\begin{itemize}[nosep]
    \item \textbf{Matemáticas (M1, M2)}: álgebra lineal y cálculo
    matricial, base teórica de la SVD.
    \item \textbf{Computación Numérica (CN)}: estabilidad numérica y
    aproximación matricial.
    \item \textbf{Inteligencia Artificial (IA)}: fundamentos de
    aprendizaje automático y evaluación de modelos.
    \item \textbf{Algorítmica (A) y Estructuras de Datos (EDA)}: diseño
    y análisis de complejidad, base del algoritmo greedy del
    Capítulo~\ref{sec:sintesis}.
    \item \textbf{Paralelismo (PAR)}: ejecución eficiente en GPU.
    \item \textbf{Projectes de Programació (PROP)}: metodologías de
    desarrollo de software y gestión incremental.
    \item \textbf{Intercambio en DTU (Dinamarca, Q1 2024--25)}:
    asignaturas de nivel máster en \emph{Machine Learning},
    \emph{Deep Learning} y \emph{Graph Neural Networks}, base técnica
    directa para el trabajo con Transformers e interpretabilidad.
\end{itemize}

Adicionalmente, durante el Q2 2025--26 se cursa simultáneamente la
asignatura \textbf{Compresión de Datos e Imagen (CDI)}, complementaria
al tema del TFG.

\subsubsection{Formulación del problema}

\begin{quote}
    \emph{¿Cómo afecta la compresión SVD selectiva ---aplicada por tipo
    de componente y por profundidad de capa--- al rendimiento y a las
    representaciones internas de un BERT fine-tuneado para
    clasificación multi-etiqueta de emociones? ¿Y qué revelan las
    técnicas de interpretabilidad mecánica sobre la organización
    interna de la información emocional?}
\end{quote}

\subsubsection{Actores implicados}

\begin{itemize}[nosep]
    \item \textbf{Investigadores e ingenieros de NLP}: beneficiarios
    primarios de los mapas de sensibilidad y del algoritmo greedy.
    \item \textbf{Comunidad académica y de código abierto}: el código
    modular basado en Hugging Face Transformers~\cite{wolf2020transformers}
    permite reutilizar la metodología para otros modelos y tareas.
    \item \textbf{Empresas que despliegan modelos de lenguaje}:
    organizaciones que integran modelos en sus productos (moderación
    de contenido, análisis de opinión, clasificación de tickets) y
    buscan reducir costes de infraestructura.
\end{itemize}


% -------------------------------------------------------------
\subsection{Justificación del enfoque}
\label{sec:gep-justificacion}
% -------------------------------------------------------------

La compresión de modelos Transformer cuenta con cinco familias
principales de técnicas: poda, cuantización, destilación, factorización
de bajo rango (SVD) y adaptación de bajo rango (LoRA). El
Capítulo~\ref{sec:estado-arte} revisa cada una con detalle. De entre
ellas, la SVD encaja con los objetivos del proyecto por cuatro razones:
(1)~no requiere re-entrenamiento, lo que permite evaluar decenas de
configuraciones a coste bajo; (2)~ofrece granularidad máxima por capa y
componente; (3)~el teorema de Eckart-Young~\cite{eckart1936approximation}
garantiza optimalidad respecto a la norma de Frobenius; y (4)~el
espectro de valores singulares aporta interpretabilidad directa.

\subsubsection{Motivación desde la sostenibilidad}

Strubell et al.~\cite{strubell2019energy} estimaron que entrenar ciertos
modelos de NLP puede emitir cantidades de CO$_2$ comparables a las de
varios coches a lo largo de su vida útil; Schwartz et
al.~\cite{schwartz2020green} acuñaron \emph{Green AI} para abogar por
priorizar la eficiencia. La compresión SVD post-entrenamiento es
afín a esta filosofía: reduce el coste de inferencia ---donde se
concentra el gasto energético a largo plazo--- sin requerir
entrenamiento adicional.


% -------------------------------------------------------------
\subsection{Alcance: objetivos y requisitos}
\label{sec:gep-alcance}
% -------------------------------------------------------------

\subsubsection{Objetivo general}

\begin{quote}
    Caracterizar cómo la compresión SVD selectiva afecta a las
    representaciones internas y al rendimiento de un BERT fine-tuneado
    para clasificación de emociones, combinar análisis de sensibilidad
    con interpretabilidad mecánica, y derivar una estrategia de
    compresión optimizada que maximice la reducción de parámetros con
    mínima pérdida de rendimiento.
\end{quote}

\subsubsection{Sub-objetivos}

\begin{enumerate}[nosep]
    \item \textbf{O1 --- Baseline competitivo}: Demostrar que
    BERT-base-uncased alcanza clasificación multi-etiqueta funcional
    sobre las 23~emociones de GoEmotions tras filtrar las 5~que
    producen $F1=0$ o no son emociones (\emph{neutral}, \emph{grief},
    \emph{nervousness}, \emph{pride}, \emph{relief}; ver
    Cap.~\ref{sec:metodologia-tfg}).
    \item \textbf{O2 --- Módulo SVD reutilizable}: Implementar
    \texttt{SVDLinear} con soporte para compresión uniforme y selectiva
    por componente, capa o ambos.
    \item \textbf{O3 --- Trade-off compresión--rendimiento}: Cuantificar
    F1 macro frente a rangos uniformes $\{512, 384, 256, 128, 64\}$.
    \item \textbf{O4 --- Visualización de la degradación}: Producir
    proyecciones t-SNE de embeddings \texttt{[CLS]}.
    \item \textbf{O5 --- Sensibilidad por componente}: Revelar qué
    tipos de capas (Q, K, V, Output, FFN-Inter, FFN-Out) son más
    vulnerables.
    \item \textbf{O6 --- Sensibilidad por profundidad}: Identificar si
    capas tempranas, medias o tardías son más críticas.
    \item \textbf{O7 --- Compresión adaptativa por energía espectral}:
    Asignar rangos no uniformes según el perfil de energía y evaluar
    en la frontera de Pareto.
    \item \textbf{O8 --- Localizar la cristalización emocional}:
    Mediante probing lineal capa a capa, determinar a qué profundidad
    cada emoción se vuelve linealmente separable. Se entrenan
    23~emociones~$\times$~12~capas~=~276~probes activos (los probes
    sobre la capa de embedding producen $F1=0$ uniformemente y se
    descartan; ver Cap.~\ref{sec:resultados-interpretabilidad}).
    \item \textbf{O9 --- Evidencia funcional de localización}: Usar
    activation patching y logit lens para distinguir información
    \emph{presente}, información \emph{leíble} por el clasificador y
    componentes \emph{suficientes} para la recuperación.
    \item \textbf{O10 --- Mapear especialización a nivel de cabeza y
    neurona}: Ablación sistemática de las 144~cabezas y análisis de
    selectividad de las 36\,864 neuronas FFN.
    \item \textbf{O11 --- Compresión informada}: Integrar O5--O10 en
    una asignación de rangos por capa y componente que supere a las
    estrategias ciegas en la frontera de Pareto.
    \item \textbf{O12 --- Recuperación por fine-tuning post-SVD}:
    Evaluar si tres épocas de re-entrenamiento sobre el modelo
    comprimido recuperan o superan al baseline original
    (Cap.~\ref{sec:sintesis}). \emph{Originalmente clasificado como
    objetivo opcional; promovido a objetivo principal tras observar el
    efecto regularizador.}
\end{enumerate}

\subsubsection{Requisitos funcionales}

\begin{itemize}[nosep]
    \item \textbf{RF1}: Compresión SVD aplicable a cualquier subconjunto
    de capas lineales del encoder.
    \item \textbf{RF2}: Evaluación automática con F1 macro y micro.
    \item \textbf{RF3}: Extracción de embeddings \texttt{[CLS]} y
    proyecciones t-SNE conjuntas.
    \item \textbf{RF4}: Cálculo de silhouette score sobre las
    proyecciones.
    \item \textbf{RF5}: Extracción de hidden states de todas las capas
    para probing y logit lens.
    \item \textbf{RF6}: Mecanismo de activation patching por capa y
    componente.
    \item \textbf{RF7}: Ablación individual de las 144 cabezas con
    medición de F1 por emoción.
    \item \textbf{RF8}: Generación de figuras y tablas exportables.
    \item \textbf{RF9}: Rangos diferenciados por componente y capa.
    \item \textbf{RF10}: Fine-tuning parcial del modelo comprimido.
\end{itemize}

\subsubsection{Requisitos no funcionales}

\begin{itemize}[nosep]
    \item \textbf{RNF1 (Reproducibilidad)}: Semilla fija (42),
    versiones de dependencias documentadas.
    \item \textbf{RNF2 (Modularidad)}: Código en módulos (\texttt{src/data},
    \texttt{src/models}, \texttt{src/compression}, \texttt{src/utils})
    separados de los notebooks.
    \item \textbf{RNF3 (Compatibilidad)}: Notebooks ejecutables en local
    y en Google Colab.
    \item \textbf{RNF4 (Eficiencia de memoria)}: Compresión sobre copias
    del modelo, liberando memoria tras cada evaluación.
    \item \textbf{RNF5 (Estabilidad numérica)}: SVD computada en
    \texttt{float32}.
\end{itemize}

\subsubsection{Obstáculos y riesgos}

\begin{table}[H]
    \centering
    \caption{Principales obstáculos y riesgos del proyecto.}
    \label{tab:gep-riesgos}
    \small
    \begin{tabular}{@{}p{3.8cm} c c p{5.5cm}@{}}
        \toprule
        \textbf{Riesgo} & \textbf{Prob.} & \textbf{Imp.} & \textbf{Mitigación} \\
        \midrule
        R1 --- Recursos computacionales insuficientes (GPU)
        & Media & Alto
        & Google Colab Pro; subsampleo para t-SNE; evaluaciones en CPU si es necesario. \\[0.2cm]
        R2 --- Inestabilidad numérica con rangos muy bajos
        & Baja & Medio
        & SVD en \texttt{float32}; validación dimensional pre-evaluación. \\[0.2cm]
        R3 --- Tiempo de ejecución excesivo
        & Alta & Medio
        & Ejecución secuencial con liberación de memoria; priorización. \\[0.2cm]
        R4 --- Métricas no capturan degradación real
        & Baja & Alto
        & F1 + silhouette + visualización cualitativa. \\[0.2cm]
        R5 --- Bug en selección de capas
        & Media & Medio
        & Tests del módulo de filtrado; alternativas documentadas. \\[0.2cm]
        R6 --- Fine-tuning post-SVD no recupera rendimiento
        & Media & Medio
        & Documentar como hallazgo válido; explorar varias configuraciones. \\[0.2cm]
        R7 --- Desviación del calendario por complejidad imprevista
        & Media & Alto
        & O1--O11 como núcleo; O12 y benchmarks como opcionales. \\
        \bottomrule
    \end{tabular}
    \\[0.15cm]{\small Fuente: elaboración propia.}
\end{table}


% -------------------------------------------------------------
\subsection{Metodología y rigor}
\label{sec:gep-metodologia}
% -------------------------------------------------------------

Dado el carácter experimental e iterativo del proyecto, se adopta una
metodología \emph{incremental iterativa}, tomando elementos de marcos
ágiles como Scrum ---estudiados en PROP--- adaptados a un contexto
unipersonal. Cada fase sigue el ciclo:
\emph{diseño}~$\to$~\emph{implementación}~$\to$~\emph{ejecución}~$\to$~%
\emph{análisis}~$\to$~\emph{validación}, con reuniones quincenales con
el director del TFG como puntos formales de revisión.

Esta elección se justifica por la naturaleza exploratoria del proyecto:
los resultados de compresión uniforme (G3) motivaron el análisis
desagregado por componente (G4); los estudios de lesión (G5) revelaron
la necesidad de técnicas causales como activation patching (G6); y la
brecha entre probing y patching motivó la incorporación del logit lens
(no prevista inicialmente) como cuarta técnica de interpretabilidad.

\subsubsection{Herramientas}

\begin{itemize}[nosep]
    \item \textbf{Control de versiones}: Git + GitHub.
    \item \textbf{Seguimiento}: GitHub Issues como backlog.
    \item \textbf{Ecosistema}: Python 3.10+, PyTorch, Hugging Face
    Transformers~\cite{wolf2020transformers}, scikit-learn,
    matplotlib, seaborn.
    \item \textbf{Dependencias fijadas}: \texttt{requirements.txt}.
    \item \textbf{Documentación}: Overleaf (\LaTeX{}).
    \item \textbf{Cómputo}: PC personal + Google Colab (GPU Tesla T4).
\end{itemize}

\subsubsection{Métodos de validación}

\begin{itemize}[nosep]
    \item \textbf{Cuantitativa}: F1 macro y micro sobre el test
    completo + ratio de parámetros.
    \item \textbf{Representacional}: silhouette score sobre
    proyecciones t-SNE.
    \item \textbf{Visual}: scatter + KDE para inspección cualitativa.
    \item \textbf{Funcional/causal}: activation patching, logit lens
    y ablación de cabezas para distinguir computación, alineación y
    transmisión.
    \item \textbf{Reproducibilidad}: semilla fija (42), código modular,
    notebooks autocontenidos.
\end{itemize}


% -------------------------------------------------------------
\subsection{Planificación temporal}
\label{sec:gep-planificacion}
% -------------------------------------------------------------

\subsubsection{Datos globales}

\begin{table}[H]
\centering
\caption{Parámetros temporales del proyecto.}
\label{tab:gep-globales}
\small
\begin{tabular}{@{}ll@{\quad}ll@{}}
\toprule
\textbf{Parámetro}          & \textbf{Valor}        & \textbf{Parámetro}                  & \textbf{Valor} \\
\midrule
Fecha de inicio              & 2~feb.\ 2026          & Horas diarias de trabajo            & 4~h/día (lectivos) \\
Fecha de fin prevista        & 26~jun.\ 2026         & Total horas núcleo                  & 457~h \\
Duración total               & 22~semanas            & Total horas opcionales              & 40~h \\
Fecha prevista de lectura    & Julio~2026            & \textbf{Total general}              & \textbf{497~h ($\approx$20 ECTS)} \\
\bottomrule
\end{tabular}
\\[0.15cm]{\small Fuente: elaboración propia.}
\end{table}

Aunque todas las tareas son ejecutadas por el autor (Guido Biosca
Lasa), se asigna a cada tarea el \emph{perfil profesional} que mejor
se ajusta al tipo de trabajo: Jefe de Proyecto (JP), Científico de
Datos (CD) o Desarrollador (Dev). Esta diferenciación refleja las
competencias requeridas y permite una estimación presupuestaria
realista (\S\ref{sec:gep-presupuesto}).

\subsubsection{Descripción de las fases}

\paragraph{G0 --- Gestión del proyecto (130~h, concurrente).}
T01--T03 (50~h, JP): tres entregas GEP. T04 (20~h, JP): seguimiento
con director (10 reuniones quincenales). T05 (60~h, JP): redacción
progresiva de la memoria.

\paragraph{G1 --- Configuración y baseline (39~h).}
T06 (8~h, Dev): repositorio + entorno. T07 (25~h, CD): fine-tuning de
BERT-base-uncased sobre GoEmotions, con filtrado a 23~emociones tras
descartar \emph{neutral} (no es emoción) y las cuatro de baja
frecuencia con $F1=0$. T08 (6~h, CD): evaluación baseline (F1 macro
0,577).

\paragraph{G2 --- Módulo de compresión SVD (30~h).}
T09 (20~h, Dev): \texttt{SVDLinear} y \texttt{apply\_svd\_compression()}
en \texttt{src/compression}. T10 (10~h, Dev): tests de validación.

\paragraph{G3 --- Compresión uniforme y análisis espectral (50~h).}
T11 (15~h, CD): experimentos uniformes en 5~rangos. T12 (10~h, CD):
perfil de energía espectral de las 72 matrices. T13 (10~h, CD):
estrategia adaptativa con 5 umbrales (80\%, 85\%, 90\%, 95\%, 99\%).
T14 (15~h, CD): frontera de Pareto y vulnerabilidad por emoción.

\paragraph{G4 --- Análisis por componente y profundidad (45~h).}
T15 (20~h, CD): 24 configuraciones (8~componentes $\times$ 3~rangos)
con extracción de embeddings. T16 (15~h, CD): 9 configuraciones
(3~bandas de profundidad $\times$ 3~rangos). T17 (10~h, CD): heatmaps,
t-SNE conjuntos, silhouette.

\paragraph{G5 --- Lesion study (23~h).}
T18 (15~h, CD): 12 lesiones por capa + 6 por componente. T19 (8~h, CD):
mapa emocional, clustering jerárquico, escáner por emoción.

\paragraph{G6 --- Interpretabilidad mecánica (65~h).}
T20 (15~h, CD): probing lineal sobre 13 hidden states; en términos
efectivos quedan 23~emociones~$\times$~12~capas~=~276~probes activos
tras descartar el embedding (F1 nulo uniforme). T21 (20~h, CD):
activation patching causal a granularidad capa y componente. T22
(10~h, CD): ablación de las 144 cabezas con taxonomía
2$\times$2 (importancia, selectividad). T23 (10~h, CD): selectividad de
las 36\,864 neuronas FFN. \emph{T20bis (10~h, CD): logit lens ---
incorporación no prevista en la planificación inicial; necesaria para
distinguir información presente, leíble y suficiente
(\S\ref{sec:logit-lens}).}

\paragraph{G7 --- Síntesis informada (55~h).}
T24 (20~h, CD): estrategia heurística informada (light, moderate,
aggressive). T25 (10~h, CD): algoritmo greedy data-driven a partir de
las sensibilidades empíricas + comparativa Pareto. T26 (25~h, CD):
fine-tuning post-SVD del modelo \texttt{greedy\_90} (3~epochs);
\emph{originalmente opcional, promovido a tarea principal de G7 tras
observar el efecto regularizador documentado en
\S\ref{sec:finetuning-recovery}}.

\paragraph{G8 --- Objetivos secundarios opcionales (15~h, no
ejecutados).}
T27 (15~h, CD, opcional): benchmarks de latencia y throughput. No
ejecutado en este TFG; documentado como trabajo futuro
(\S\ref{sec:trabajo-futuro}).

\paragraph{G9 --- Cierre y síntesis (55~h).}
T28 (40~h, JP): integración final de la memoria. T29 (15~h, JP):
preparación de la presentación oral.

\subsubsection{Tabla resumen de tareas}

\begin{longtable}{@{}p{0.55cm} p{4.6cm} c c p{2.3cm} p{1.5cm}@{}}
    \caption{Tabla resumen de tareas.}
    \label{tab:gep-tareas}\\
    \toprule
    \textbf{Cód.} & \textbf{Tarea} & \textbf{Perfil} & \textbf{h} & \textbf{Dependencias} & \textbf{Rec.\ mat.} \\
    \midrule
    \endfirsthead
    \multicolumn{6}{c}{\footnotesize(Tabla~\ref{tab:gep-tareas}, continuación)}\\
    \toprule
    \textbf{Cód.} & \textbf{Tarea} & \textbf{Perfil} & \textbf{h} & \textbf{Dependencias} & \textbf{Rec.\ mat.} \\
    \midrule
    \endhead
    \midrule\endfoot
    \bottomrule\\[-4pt]
    \multicolumn{6}{@{}l}{\footnotesize [opt]=opcional; [conc.]=concurrente. Rec.\ humano en todas: Guido Biosca Lasa. Fuente: elab.\ propia.}
    \endlastfoot
    \rowcolor{ganttG0!40}\multicolumn{6}{@{}l}{\footnotesize\textbf{G0 --- Gestión (130~h)}}\\
    T01 & Contextualización (E1)               & JP  & 20 & ---     & Overleaf \\
    T02 & Planificación (E2)                    & JP  & 15 & T01     & Overleaf \\
    T03 & Presupuesto + sostenibilidad (E3)     & JP  & 15 & T01     & Overleaf \\
    T04 & Reuniones director [conc.]            & JP  & 20 & T01     & Videoconf. \\
    T05 & Redacción memoria [conc.]             & JP  & 60 & T01     & Overleaf \\
    \midrule
    \rowcolor{ganttG1!40}\multicolumn{6}{@{}l}{\footnotesize\textbf{G1 --- Baseline (39~h)}}\\
    T06 & Config.\ entorno + repo               & Dev &  8 & T01     & PC, GitHub \\
    T07 & Fine-tuning BERT (23 emociones)       & CD  & 25 & T06     & GPU \\
    T08 & Evaluación baseline                   & CD  &  6 & T07     & PC \\
    \midrule
    \rowcolor{ganttG2!40}\multicolumn{6}{@{}l}{\footnotesize\textbf{G2 --- Módulo SVD (30~h)}}\\
    T09 & Impl.\ \texttt{SVDLinear}              & Dev & 20 & T08     & PC \\
    T10 & Tests validación                      & Dev & 10 & T09     & PC \\
    \midrule
    \rowcolor{ganttGE!30}\multicolumn{6}{@{}l}{\footnotesize\textbf{G3 --- Uniforme + espectral (50~h)}}\\
    T11 & Comp.\ uniforme (5 rangos)            & CD  & 15 & T10     & GPU \\
    T12 & Análisis espectral 72 matrices        & CD  & 10 & T10     & PC \\
    T13 & Adaptativa (5 umbrales energéticos)   & CD  & 10 & T12     & PC \\
    T14 & Pareto + vulnerabilidad emocional     & CD  & 15 & T11,T13 & GPU \\
    \midrule
    \rowcolor{ganttGE!45}\multicolumn{6}{@{}l}{\footnotesize\textbf{G4 --- Comp.\ + profundidad (45~h)}}\\
    T15 & Componente (24 configs)               & CD  & 20 & T11     & GPU \\
    T16 & Profundidad (9 configs)               & CD  & 15 & T15     & GPU \\
    T17 & Heatmaps + t-SNE + silhouette         & CD  & 10 & T16     & PC \\
    \midrule
    \rowcolor{ganttGE!55}\multicolumn{6}{@{}l}{\footnotesize\textbf{G5 --- Lesion study (23~h)}}\\
    T18 & Lesiones capa (12) + componente (6)   & CD  & 15 & T10     & GPU \\
    T19 & Mapa emocional + clustering           & CD  &  8 & T18     & PC \\
    \midrule
    \rowcolor{ganttGE!70}\multicolumn{6}{@{}l}{\footnotesize\textbf{G6 --- Interpretabilidad (65~h)}}\\
    T20 & Probing lineal (276 probes activos)    & CD  & 15 & T08     & PC \\
    T20bis & Logit lens [no previsto inicialmente] & CD & 10 & T20  & PC \\
    T21 & Activation patching                    & CD  & 20 & T11     & GPU \\
    T22 & Ablación cabezas (144)                 & CD  & 10 & T08     & GPU \\
    T23 & Análisis neuronal FFN (36\,864)        & CD  & 10 & T08     & GPU \\
    \midrule
    \rowcolor{ganttGE!85}\multicolumn{6}{@{}l}{\footnotesize\textbf{G7 --- Síntesis informada (55~h)}}\\
    T24 & Heurística informada (3 variantes)     & CD  & 20 & T17,T19--T23 & GPU \\
    T25 & Greedy data-driven + Pareto comp.      & CD  & 10 & T24     & PC \\
    T26 & Fine-tuning post-SVD [promov.\ a obligatoria] & CD & 25 & T25 & GPU \\
    \midrule
    \rowcolor{ganttGO!50}\multicolumn{6}{@{}l}{\footnotesize\textbf{G8 --- Opcional no ejecutado (15~h)}}\\
    T27 & Benchmarks latencia/throughput [opt]   & CD  & 15 & T25     & PC+GPU \\
    \midrule
    \rowcolor{ganttG9!50}\multicolumn{6}{@{}l}{\footnotesize\textbf{G9 --- Cierre (55~h)}}\\
    T28 & Integración + revisión memoria         & JP  & 40 & T26     & Overleaf \\
    T29 & Preparación presentación oral          & JP  & 15 & T28     & PC \\
    \midrule
    \multicolumn{3}{@{}l}{\textbf{Total núcleo (T01--T26, T28--T29)}} & \textbf{492} & & \\
    \multicolumn{3}{@{}l}{\textbf{Total opcional (T27)}} & \textbf{15} & & \\
    \multicolumn{3}{@{}l}{\textbf{Total general planificado}} & \textbf{507} & & \\
\end{longtable}

\subsubsection{Gestión del riesgo: planes alternativos}

Los planes B se activan únicamente si el riesgo se materializa
(Tabla~\ref{tab:gep-planesb}).

\begin{table}[H]
\centering
\caption{Planes alternativos asociados a los riesgos R1--R7.}
\label{tab:gep-planesb}
\footnotesize
\begin{tabular}{@{}p{2.1cm} c c p{5.3cm} p{1.5cm} p{1.2cm}@{}}
\toprule
\textbf{Riesgo} & \textbf{Prob.} & \textbf{Imp.} & \textbf{Plan B} & \textbf{Rec.\ adic.} & \textbf{$\Delta$h} \\
\midrule
R1 GPU insuf.       & Media & Alto  & Reducir T15 a 4 componentes; subsamplear t-SNE a 1\,500 puntos. & Ninguno & $-$13~h \\
\midrule
R2 Inest.\ num.     & Baja  & Medio & Validación numérica adicional (5~h) en \texttt{float64}. & 5~h & $+$5~h \\
\midrule
R3 Tiempo excesivo  & Alta  & Medio & Reducir T15--T16; paralelizar T11/T18. & Ninguno & $-$15~h \\
\midrule
R4 Métricas insuf.  & Baja  & Alto  & Análisis cualitativo (8~h): inspección de errores. & 8~h & $+$8~h \\
\midrule
R5 Bug capas        & Media & Medio & T10 incluye tests; indexación manual de respaldo (4~h). & 4~h & $+$4~h \\
\midrule
R6 FT no recupera   & Media & Medio & Documentar como resultado válido. T26 originalmente opcional. & Ninguno & 0~h \\
\midrule
R7 Desviación gral. & Media & Alto  & Eliminar T27 ($-$15~h); reducir T22+T23; documentar como trabajo futuro. & Ninguno & $-$15--40~h \\
\bottomrule
\end{tabular}
\\[0.15cm]{\small En el peor escenario (R1+R3+R7), el núcleo se reduce a $\approx$370~h, garantizando O1--O9. Fuente: elab.\ propia.}
\end{table}

\subsubsection{Diagrama de Gantt}

La cadena de dependencias crítica es:
\[
\text{T06} \to \text{T07} \to \text{T08} \to \text{T09} \to \text{T10}
\to \text{T11} \to \text{T15} \to \text{T16} \to \text{T17} \to \text{T24} \to
\text{T25} \to \text{T26} \to \text{T28} \to \text{T29}.
\]
T04+T05 (gestión, redacción) son concurrentes durante todas las
22~semanas. T12 ejecuta en paralelo con T11. G5 y parte de G6 son
ejecutables en paralelo con G3--G4.

\begin{landscape}
\begin{figure}[H]
\centering
\resizebox{\linewidth}{!}{%
\begin{ganttchart}[
    hgrid, vgrid,
    x unit=0.9cm, y unit title=0.62cm, y unit chart=0.44cm,
    title height=1,
    title label font=\scriptsize\bfseries,
    bar label font=\scriptsize,
    group label font=\bfseries\scriptsize,
    milestone label font=\itshape\scriptsize,
    bar height=0.65, bar top shift=0.175,
    group right shift=0, group top shift=0.65, group height=0.35, group peaks height=0.2,
    milestone height=0.55, milestone top shift=0.225,
]{1}{22}
\gantttitle{Febrero}{4}\gantttitle{Marzo}{4}\gantttitle{Abril}{5}\gantttitle{Mayo}{4}\gantttitle{Junio}{5}\\
\gantttitlelist{1,...,22}{1}\\
%--- G0 ---
\ganttgroup[group/.append style={fill=ganttG0,draw=black}]{G0: Gestión}{1}{22}\\
\ganttbar[name=t01,bar/.append style={fill=ganttG0}]{T01--T03 Entregas GEP}{1}{7}\\
\ganttbar[name=t04,bar/.append style={fill=ganttG0!60,draw=gray!60}]{T04+T05 Reuniones+Doc.\ [conc.]}{1}{22}\\
%--- G1 ---
\ganttgroup[group/.append style={fill=ganttG1,draw=black}]{G1: Baseline}{1}{7}\\
\ganttbar[name=t06,bar/.append style={fill=ganttG1}]{T06 Config.\ entorno}{1}{2}\\
\ganttbar[name=t07,bar/.append style={fill=ganttG1}]{T07 Fine-tuning BERT}{2}{6}\\
\ganttbar[name=t08,bar/.append style={fill=ganttG1}]{T08 Eval.\ baseline}{6}{7}\\
%--- G2 ---
\ganttgroup[group/.append style={fill=ganttG2,draw=black}]{G2: Módulo SVD}{7}{10}\\
\ganttbar[name=t09,bar/.append style={fill=ganttG2}]{T09 Impl.\ SVDLinear}{7}{9}\\
\ganttbar[name=t10,bar/.append style={fill=ganttG2}]{T10 Tests validación}{9}{10}\\
%--- G3 ---
\ganttgroup[group/.append style={fill=ganttGE!70,draw=black}]{G3: Uniforme+Espectral}{10}{13}\\
\ganttbar[name=t11,bar/.append style={fill=ganttGE!70}]{T11 Comp.\ uniforme}{10}{12}\\
\ganttbar[name=t12,bar/.append style={fill=ganttGE!50}]{T12 Análisis espectral}{10}{11}\\
\ganttbar[name=t13,bar/.append style={fill=ganttGE!50}]{T13--T14 Adaptativa+Pareto}{12}{13}\\
%--- G4 ---
\ganttgroup[group/.append style={fill=ganttGE,draw=black}]{G4: Comp.+Prof.}{12}{16}\\
\ganttbar[name=t15,bar/.append style={fill=ganttGE}]{T15 Exp.\ componente (24)}{12}{14}\\
\ganttbar[name=t16,bar/.append style={fill=ganttGE}]{T16 Exp.\ profundidad (9)}{14}{15}\\
\ganttbar[name=t17,bar/.append style={fill=ganttGE}]{T17 Heatmaps+t-SNE}{15}{16}\\
%--- G5 ---
\ganttgroup[group/.append style={fill=ganttGE!85,draw=black}]{G5: Lesion study}{10}{13}\\
\ganttbar[name=t18,bar/.append style={fill=ganttGE!85}]{T18--T19 Lesiones+Mapa}{10}{13}\\
%--- G6 ---
\ganttgroup[group/.append style={fill=ganttGE!90,draw=black}]{G6: Interpretabilidad}{8}{17}\\
\ganttbar[name=t20,bar/.append style={fill=ganttGE!90}]{T20 Probing lineal}{8}{10}\\
\ganttbar[name=t20b,bar/.append style={fill=ganttGE!95}]{T20bis Logit lens [nuevo]}{14}{15}\\
\ganttbar[name=t21,bar/.append style={fill=ganttGE!90}]{T21 Act.\ patching}{12}{14}\\
\ganttbar[name=t22,bar/.append style={fill=ganttGE!90}]{T22--T23 Heads+Neurons}{10}{13}\\
%--- G7 ---
\ganttgroup[group/.append style={fill=ganttGE,draw=black}]{G7: Síntesis}{16}{20}\\
\ganttbar[name=t24,bar/.append style={fill=ganttGE}]{T24--T25 Informada+Greedy}{16}{18}\\
\ganttbar[name=t26,bar/.append style={fill=ganttGE}]{T26 Fine-tuning post-SVD}{18}{20}\\
%--- G9 ---
\ganttgroup[group/.append style={fill=ganttG9,draw=black}]{G9: Cierre}{20}{22}\\
\ganttbar[name=t28,bar/.append style={fill=ganttG9}]{T28 Memoria final}{20}{22}\\
\ganttbar[name=t29,bar/.append style={fill=ganttG9}]{T29 Presentación oral}{22}{22}\\
%--- Hito ---
\ganttmilestone[name=m1,milestone/.append style={fill=red!80,draw=red}]{Lectura TFG}{22}\\
%--- Dependencias ---
\ganttlink{t06}{t07}\ganttlink{t07}{t08}\ganttlink{t08}{t09}\ganttlink{t09}{t10}
\ganttlink{t10}{t11}\ganttlink{t11}{t15}\ganttlink{t15}{t16}\ganttlink{t16}{t17}
\ganttlink{t17}{t24}\ganttlink{t24}{t26}\ganttlink{t26}{t28}\ganttlink{t28}{t29}
\end{ganttchart}
}%
\caption{Diagrama de Gantt actualizado. T20bis (logit lens) y T26 (fine-tuning) reflejan los cambios respecto a la planificación inicial. Grises: opcionales no ejecutados. Flechas: cadena crítica.}
\label{fig:gep-gantt}
\end{figure}
\end{landscape}


% -------------------------------------------------------------
\subsection{Gestión económica (presupuesto)}
\label{sec:gep-presupuesto}
% -------------------------------------------------------------

El presupuesto se estructura en cuatro partidas siguiendo la
metodología del módulo: costes de personal por actividad (CPA),
costes genéricos (CG), contingencia e imprevistos.

\subsubsection{Costes por hora por perfil profesional}

\begin{table}[H]
\centering
\caption{Coste por hora estimado por perfil profesional.}
\label{tab:gep-perfiles}
\small
\begin{tabular}{@{}l r r r@{}}
\toprule
\textbf{Perfil} & \textbf{Bruto anual} & \textbf{Con SS ($\times$1,3)} & \textbf{Coste/hora} \\
\midrule
Jefe de proyecto (JP)    & 32.000\,\euro{} & 41.600\,\euro{} & 24,19\,\euro{}/h $\approx$ \textbf{24\,\euro{}/h} \\
Científico de datos (CD) & 28.000\,\euro{} & 36.400\,\euro{} & 21,16\,\euro{}/h $\approx$ \textbf{21\,\euro{}/h} \\
Desarrollador (Dev)      & 25.000\,\euro{} & 32.500\,\euro{} & 18,90\,\euro{}/h $\approx$ \textbf{19\,\euro{}/h} \\
\bottomrule
\end{tabular}
\\[0.15cm]{\small Fuente: elaboración propia a partir de Hays España 2025 y Glassdoor España. 1.720~horas laborales anuales.}
\end{table}

\subsubsection{Costes de personal por actividad (CPA)}

\begin{longtable}{@{}p{0.5cm} p{4.2cm} c c r r@{}}
    \caption{Costes de personal por actividad (CPA).}
    \label{tab:gep-cpa}\\
    \toprule
    \textbf{Cód.} & \textbf{Tarea} & \textbf{Perfil} & \textbf{h} & \textbf{\euro{}/h} & \textbf{Coste} \\
    \midrule
    \endfirsthead
    \multicolumn{6}{c}{\footnotesize(Tabla~\ref{tab:gep-cpa}, continuación)}\\
    \toprule
    \textbf{Cód.} & \textbf{Tarea} & \textbf{Perfil} & \textbf{h} & \textbf{\euro{}/h} & \textbf{Coste} \\
    \midrule
    \endhead
    \midrule\endfoot
    \bottomrule
    \endlastfoot
    \rowcolor{gray!10}\multicolumn{5}{@{}l}{\footnotesize\textbf{G0 --- Gestión (130~h)}} & \textbf{3.120\,\euro{}} \\
    T01 & Contextualización                  & JP  & 20 & 24 & 480\,\euro{} \\
    T02 & Planificación                       & JP  & 15 & 24 & 360\,\euro{} \\
    T03 & Presupuesto + sostenibilidad        & JP  & 15 & 24 & 360\,\euro{} \\
    T04 & Reuniones director                 & JP  & 20 & 24 & 480\,\euro{} \\
    T05 & Redacción memoria                  & JP  & 60 & 24 & 1.440\,\euro{} \\
    \midrule
    \rowcolor{gray!10}\multicolumn{5}{@{}l}{\footnotesize\textbf{G1 --- Baseline (39~h)}} & \textbf{803\,\euro{}} \\
    T06 & Config.\ entorno                    & Dev &  8 & 19 & 152\,\euro{} \\
    T07 & Fine-tuning BERT (23 emociones)     & CD  & 25 & 21 & 525\,\euro{} \\
    T08 & Eval.\ baseline                     & CD  &  6 & 21 & 126\,\euro{} \\
    \midrule
    \rowcolor{gray!10}\multicolumn{5}{@{}l}{\footnotesize\textbf{G2 --- Módulo SVD (30~h)}} & \textbf{570\,\euro{}} \\
    T09 & Impl.\ \texttt{SVDLinear}            & Dev & 20 & 19 & 380\,\euro{} \\
    T10 & Tests validación                    & Dev & 10 & 19 & 190\,\euro{} \\
    \midrule
    \rowcolor{gray!10}\multicolumn{5}{@{}l}{\footnotesize\textbf{G3 --- Uniforme + espectral (50~h)}} & \textbf{1.050\,\euro{}} \\
    T11--T14 & Uniforme, espectro, adapt., Pareto & CD & 50 & 21 & 1.050\,\euro{} \\
    \midrule
    \rowcolor{gray!10}\multicolumn{5}{@{}l}{\footnotesize\textbf{G4 --- Comp.\ + profundidad (45~h)}} & \textbf{945\,\euro{}} \\
    T15--T17 & Componente, profundidad, t-SNE  & CD & 45 & 21 & 945\,\euro{} \\
    \midrule
    \rowcolor{gray!10}\multicolumn{5}{@{}l}{\footnotesize\textbf{G5 --- Lesion study (23~h)}} & \textbf{483\,\euro{}} \\
    T18--T19 & Lesiones + mapa emocional        & CD & 23 & 21 & 483\,\euro{} \\
    \midrule
    \rowcolor{gray!10}\multicolumn{5}{@{}l}{\footnotesize\textbf{G6 --- Interpretabilidad (65~h)}} & \textbf{1.365\,\euro{}} \\
    T20 & Probing lineal (276 probes)          & CD & 15 & 21 & 315\,\euro{} \\
    T20bis & Logit lens [no previsto]          & CD & 10 & 21 & 210\,\euro{} \\
    T21 & Activation patching                  & CD & 20 & 21 & 420\,\euro{} \\
    T22 & Ablación cabezas                     & CD & 10 & 21 & 210\,\euro{} \\
    T23 & Análisis neuronal FFN                & CD & 10 & 21 & 210\,\euro{} \\
    \midrule
    \rowcolor{gray!10}\multicolumn{5}{@{}l}{\footnotesize\textbf{G7 --- Síntesis informada (55~h)}} & \textbf{1.155\,\euro{}} \\
    T24 & Heurística informada                 & CD & 20 & 21 & 420\,\euro{} \\
    T25 & Greedy + Pareto comp.                & CD & 10 & 21 & 210\,\euro{} \\
    T26 & Fine-tuning post-SVD                 & CD & 25 & 21 & 525\,\euro{} \\
    \midrule
    \rowcolor{gray!10}\multicolumn{5}{@{}l}{\footnotesize\textbf{G9 --- Cierre (55~h)}} & \textbf{1.320\,\euro{}} \\
    T28 & Memoria final                        & JP & 40 & 24 & 960\,\euro{} \\
    T29 & Presentación oral                    & JP & 15 & 24 & 360\,\euro{} \\
    \midrule
    \multicolumn{3}{@{}l}{\textbf{Total CPA núcleo (sin opcional T27)}} & \textbf{492} & & \textbf{10.811\,\euro{}} \\
\end{longtable}

\subsubsection{Costes genéricos (CG)}

\begin{table}[H]
\centering
\caption{Costes genéricos del proyecto.}
\label{tab:gep-cg}
\small
\begin{tabular}{@{}l l r r r@{}}
\toprule
\textbf{Concepto} & \textbf{Detalle} & \textbf{Coste total} & \textbf{Vida útil} & \textbf{Imputado} \\
\midrule
Amortización PC portátil & 1.200\,\euro{}, uso 5 meses & 1.200\,\euro{} & 4 años (48~m.) & 125\,\euro{} \\
Google Colab Pro         & 11,99\,\euro{}/mes $\times$ 5 meses & --- & --- & 60\,\euro{} \\
Electricidad             & 0,1\,kW $\times$ 492\,h $\times$ 0,20\,\euro{}/kWh & --- & --- & 10\,\euro{} \\
Software open source     & Python, PyTorch, HF, scikit-learn & --- & --- & 0\,\euro{} \\
Overleaf, GitHub         & Plan gratuito & --- & --- & 0\,\euro{} \\
Internet                 & Uso parcial conexión doméstica & --- & --- & 25\,\euro{} \\
Espacio de trabajo       & Uso parcial domicilio (10\,m$^2$, 5 meses) & --- & --- & 50\,\euro{} \\
\midrule
\multicolumn{4}{@{}l}{\textbf{Total costes genéricos}} & \textbf{270\,\euro{}} \\
\bottomrule
\end{tabular}
\\[0.15cm]{\small Amortización: $1.200 \times 5/48 = 125$\,\euro{}. Trabajo desde domicilio (sin desplazamientos). Fuente: elab.\ propia.}
\end{table}

\subsubsection{Contingencia e imprevistos}

Se aplica un ratio de contingencia del \textbf{15\%} sobre la suma de
CPA y CG, dentro del rango habitual del 10--20\% en proyectos
informáticos:
\[
\text{Contingencia} = (10.811 + 270) \times 0{,}15 = \textbf{1.662\,\euro{}}.
\]

Las partidas de imprevistos se vinculan a los riesgos identificados,
ponderadas por probabilidad:

\begin{table}[H]
\centering
\caption{Partidas de imprevistos.}
\label{tab:gep-imprevistos}
\small
\begin{tabular}{@{}l c r r r@{}}
\toprule
\textbf{Riesgo} & \textbf{Prob.} & \textbf{Coste si ocurre} & \textbf{Detalle} & \textbf{Imputado} \\
\midrule
R1 GPU insuf.       & 50\% & 150\,\euro{} & Upgrade Colab Pro+ (3~m.) & 75\,\euro{} \\
R2 Inest.\ num.     & 20\% & 105\,\euro{} & 5~h extra CD              & 21\,\euro{} \\
R4 Métricas insuf.  & 20\% & 168\,\euro{} & 8~h extra CD              & 34\,\euro{} \\
R5 Bug capas        & 50\% & 76\,\euro{}  & 4~h extra Dev             & 38\,\euro{} \\
\midrule
\multicolumn{4}{@{}l}{\textbf{Total imprevistos}} & \textbf{168\,\euro{}} \\
\bottomrule
\end{tabular}
\\[0.15cm]{\small R3, R6 y R7 se mitigan reduciendo alcance, sin coste adicional. Fuente: elab.\ propia.}
\end{table}

\subsubsection{Presupuesto total}

\begin{table}[H]
\centering
\caption{Resumen del presupuesto total del proyecto.}
\label{tab:gep-total}
\small
\begin{tabular}{@{}l r r@{}}
\toprule
\textbf{Partida} & \textbf{Importe} & \textbf{\%} \\
\midrule
Costes de personal (CPA)    & 10.811\,\euro{} & 84{,}1\% \\
Costes genéricos (CG)       & 270\,\euro{}    & 2{,}1\% \\
Contingencia (15\%)          & 1.662\,\euro{}  & 12{,}9\% \\
Imprevistos                  & 168\,\euro{}    & 1{,}3\% \\
\midrule
\textbf{Total presupuesto}  & \textbf{12.911\,\euro{}} & \textbf{100\%} \\
\bottomrule
\end{tabular}
\\[0.15cm]{\small El coste de personal representa $\sim$84\% del total, habitual en proyectos de investigación aplicada. Fuente: elab.\ propia.}
\end{table}

\subsubsection{Control de gestión}

Para detectar y corregir desviaciones presupuestarias durante la
ejecución se definieron los siguientes mecanismos.

\textbf{Procedimiento.} Al finalizar cada fase del Gantt (G1--G9) se
registran las horas reales dedicadas a cada tarea y se comparan con
las estimadas. Las reuniones quincenales con el director (T04) sirven
como puntos formales de revisión.

\textbf{Indicadores.}

\begin{itemize}[nosep, leftmargin=1.2em]
    \item \textbf{Desviación en horas por tarea}:
    $\Delta h_i = h_i^{\text{real}} - h_i^{\text{est}}$. Si
    $\Delta h_i > 0{,}2 \times h_i^{\text{est}}$, se activa una revisión
    de la planificación restante.
    \item \textbf{Desviación en coste por fase}: si
    $\Delta C_g > 0{,}15 \times C_g^{\text{est}}$, se evalúa el plan B
    correspondiente.
    \item \textbf{Índice de rendimiento del coste (CPI)}:
    $\text{CPI} = \text{EV}/\text{AC}$. Un CPI~$<0{,}85$ indica
    sobrecostes que requieren acción correctiva.
    \item \textbf{\% de contingencia consumida}: si supera el 50\% antes
    de completar el 75\% del proyecto, se priorizan O1--O11 y se
    eliminan tareas opcionales.
\end{itemize}


% -------------------------------------------------------------
\subsection{Sostenibilidad y compromiso social}
\label{sec:gep-sostenibilidad}
% -------------------------------------------------------------

\subsubsection{Autoevaluación}

Como \textbf{puntos fuertes}, dispongo de sensibilidad hacia el impacto
ambiental de la computación: el propio tema del TFG ---compresión de
modelos--- está motivado en parte por la reducción del coste energético
de la inferencia, y he incorporado referencias a Strubell et
al.~\cite{strubell2019energy} y Schwartz et al.~\cite{schwartz2020green}
sobre \emph{Green AI} desde la definición del proyecto. A nivel
económico, he estimado el presupuesto con detalle, diferenciando
perfiles y justificando cada partida.

Como \textbf{puntos débiles}, mi formación en sostenibilidad ha sido
más teórica que práctica: no había realizado anteriormente un análisis
de ciclo de vida energético de un proyecto informático y mi
conocimiento sobre normativas ambientales aplicables a centros de datos
es limitado. La dimensión social ---más allá de la accesibilidad---
requiere reflexión más profunda sobre sesgos en los datos de emociones
y sobre quién se beneficia realmente de estas tecnologías.

\subsubsection{Dimensión económica}

\textbf{PPP --- Coste estimado del proyecto.} El presupuesto total de
12.911\,\euro{} es razonable para un proyecto académico de
$\sim$497~h. El coste está dominado por el recurso humano (84\%), lo
que refleja un proyecto de investigación aplicada sin necesidad de
infraestructura costosa. El uso de herramientas open source y servicios
gratuitos o de bajo coste (Google Colab Pro, GitHub, Overleaf) mantiene
los costes genéricos en mínimos.

\textbf{VU --- ¿Cómo se resuelven actualmente los costes del problema?}
Desplegar modelos Transformer de gran tamaño requiere GPUs costosas
(A100, H100) con consumos energéticos elevados. Empresas y centros de
investigación invierten miles de euros mensuales en infraestructura
cloud para servir estos modelos. Las técnicas de compresión existentes
(destilación, cuantización) requieren re-entrenamiento costoso o
hardware especializado.

\textbf{¿En qué mejorará económicamente la solución?} La compresión SVD
post-entrenamiento que estudia este proyecto no requiere
re-entrenamiento ni hardware especial. El mapa de sensibilidad
producido permite a ingenieros comprimir selectivamente solo los
componentes menos críticos, maximizando la reducción de parámetros con
mínima pérdida de calidad, y por tanto reduciendo costes de inferencia
sin inversión adicional en entrenamiento.

\subsubsection{Dimensión ambiental}

\textbf{PPP --- Impacto ambiental de la realización del proyecto.} El
principal impacto proviene del consumo eléctrico de las GPUs durante
los experimentos. Estimamos un uso de GPU de $\sim$80--100~horas. Con
un consumo típico de una GPU T4 en Colab de $\sim$70\,W, esto supone
$\sim$7\,kWh. Sumando el PC personal (492~h~$\times$~0,1\,kW =
49\,kWh), el consumo total estimado es de $\sim$56\,kWh, equivalente a
$\sim$12\,kg de CO$_2$ (factor de emisión medio en España: 0,22\,kg
CO$_2$/kWh). Es un impacto modesto comparado con el entrenamiento de
modelos grandes.

\textbf{¿Te has planteado minimizar el impacto?} Sí. Se reutiliza un
modelo BERT-base ya pre-entrenado por Google, evitando el coste enorme
del pre-entrenamiento. La compresión SVD se aplica post-hoc sobre el
modelo fine-tuneado, sin requerir entrenamiento adicional para la
mayoría de experimentos. Se liberan copias del modelo de memoria tras
cada evaluación (RNF4) y se utiliza subsampleo para visualizaciones
(1.500 muestras en t-SNE).

\textbf{VU --- Mejora ambiental.} La contribución ambiental principal
del proyecto es indirecta: los mapas de sensibilidad permiten comprimir
modelos de forma informada, reduciendo el tamaño y el consumo
energético durante la inferencia. Dado que la inferencia representa la
mayor parte del gasto energético a largo plazo de un modelo desplegado,
guiar la compresión de forma óptima tiene un impacto ambiental positivo
acumulativo significativo, alineado con la filosofía
\emph{Green AI}~\cite{schwartz2020green}.

\subsubsection{Dimensión social}

\textbf{PPP --- Aportación personal.} Este proyecto me permite
profundizar en tres áreas clave para mi desarrollo profesional:
(1)~comprensión profunda de la arquitectura Transformer y sus
mecanismos internos, conocimiento fundamental en la IA actual;
(2)~habilidades en álgebra lineal aplicada y análisis numérico,
conectando teoría matemática con aplicaciones prácticas; y
(3)~experiencia en diseño y ejecución de estudios experimentales
sistemáticos.

\textbf{VU --- Mejora social.} Actualmente, los modelos de lenguaje
grandes son accesibles principalmente para grandes empresas con
recursos computacionales significativos, lo que genera una brecha
tecnológica. La compresión eficiente democratiza el acceso a estos
modelos, permitiendo su despliegue en dispositivos con recursos
limitados (móviles, sistemas embebidos) y en organizaciones con
presupuestos reducidos (universidades, ONGs, PYMES). Un modelo de
clasificación de emociones comprimido podría usarse en aplicaciones de
salud mental, moderación de contenido o análisis de opinión en
contextos donde antes era inviable.

\textbf{VU --- Necesidad real.} Sí. La tendencia hacia modelos cada
vez más grandes hace que la compresión sea no solo deseable sino
necesaria para su adopción generalizada. Además, la clasificación de
emociones tiene aplicaciones reales en salud mental, educación y
moderación de plataformas. Este proyecto no solo demuestra
\emph{cómo} comprimir, sino \emph{dónde} hacerlo de forma segura, lo
cual es una contribución con utilidad práctica directa.


% -------------------------------------------------------------
\subsection{Ejecución real frente a planificación}
\label{sec:gep-ejecucion-real}
% -------------------------------------------------------------

Esta sección, redactada al cierre del proyecto, contrasta lo
ejecutado con la planificación inicial.

\subsubsection{Ajustes de alcance respecto a la Entrega 4}

Durante la ejecución se introdujeron cuatro modificaciones de alcance
respecto a lo recogido en la Entrega 4 GEP, todas en la dirección de
\emph{ampliar la profundidad del análisis} y todas justificadas por los
hallazgos intermedios:

\begin{enumerate}[nosep]
    \item \textbf{Filtrado a 23 emociones (afecta O1, O3, T07, T08).}
    Tras el fine-tuning sobre las 28 emociones de GoEmotions, cuatro
    emociones (\emph{grief}, \emph{nervousness}, \emph{pride},
    \emph{relief}) producían $F1=0$ y la categoría \emph{neutral}
    distorsionaba las métricas macro. Se decidió filtrar las cinco y
    trabajar con 23 emociones funcionales (Cap.~\ref{sec:metodologia-tfg}).
    El filtrado no añadió horas significativas de cómputo.

    \item \textbf{Probes efectivos: 276 en lugar de 364.} La
    planificación inicial estimaba 28~emociones $\times$ 13~capas. Tras
    el filtrado a 23~emociones y el descarte del embedding (donde el
    F1 lineal es nulo uniformemente), el número efectivo de
    clasificadores entrenados es 276 (Cap.~\ref{sec:resultados-interpretabilidad}).
    El coste computacional se redujo en consecuencia.

    \item \textbf{Logit lens (T20bis, +10\,h, no prevista).} La brecha
    entre el probing (información presente) y el activation patching
    (suficiencia funcional) motivó incorporar el logit lens como cuarta
    técnica de interpretabilidad para distinguir información
    \emph{presente} de información \emph{leíble} por la cabeza de
    clasificación entrenada (\S\ref{sec:logit-lens}). Esto extendió G6
    de 55\,h a 65\,h.

    \item \textbf{Fine-tuning post-SVD (T26) promovido a tarea
    principal.} Inicialmente clasificada como opcional dentro de G8,
    el \emph{fine-tuning} de tres épocas sobre el modelo
    \texttt{greedy\_90} produjo una observación que merecía análisis:
    el modelo comprimido al 86{,}4\% de los parámetros alcanza un F1
    macro comparable al baseline original (0{,}591 vs 0{,}577), con
    una mejora sistemática y notable en las emociones
    infrarrepresentadas (\S\ref{sec:finetuning-recovery}). El alcance
    causal del hallazgo está pendiente de un experimento de control
    (baseline + 3 epochs adicionales) propuesto en
    \S\ref{sec:trabajo-futuro}. Dada la relevancia descriptiva, T26
    se reclasificó como tarea principal de G7; como contrapartida,
    T27 (benchmarks de latencia) se pospuso como trabajo futuro.
\end{enumerate}

\subsubsection{Resumen del balance de horas}

\begin{table}[H]
\centering
\caption{Comparativa entre planificación inicial (Entrega 4) y ejecución
real.}
\label{tab:gep-balance-horas}
\small
\begin{tabular}{@{}l r r r@{}}
\toprule
\textbf{Fase} & \textbf{Plan inicial (h)} & \textbf{Ejecutado (h)} & \textbf{$\Delta$} \\
\midrule
G0 Gestión                    & 130 & 130 & 0 \\
G1 Baseline                   & 39  & 39  & 0 \\
G2 Módulo SVD                 & 30  & 30  & 0 \\
G3 Uniforme + espectral       & 50  & 50  & 0 \\
G4 Componente + profundidad   & 45  & 45  & 0 \\
G5 Lesion study               & 23  & 23  & 0 \\
G6 Interpretabilidad          & 55  & 65  & $+$10 (logit lens) \\
G7 Síntesis informada         & 30  & 55  & $+$25 (T26 promovido) \\
G8 Opcionales                 & 40  & 0   & $-$40 (T27 no ejecutado) \\
G9 Cierre                     & 55  & 55  & 0 \\
\midrule
\textbf{Total}                & \textbf{497} & \textbf{492} & $-$5 \\
\bottomrule
\end{tabular}
\\[0.15cm]{\small La incorporación del logit lens y la promoción del fine-tuning post-SVD se compensaron con la exclusión de los benchmarks de latencia, manteniendo el total dentro del rango planificado.}
\end{table}

\subsubsection{Planes B activados}

De los 7 riesgos identificados, sólo \textbf{R3 (tiempo de ejecución
excesivo)} se materializó parcialmente en G4: el experimento por
componente (T15) se ejecutó secuencialmente con liberación explícita de
memoria entre configuraciones para evitar saturar la GPU. No fue
necesario reducir el alcance de T15.

\subsubsection{Indicadores de control finales}

\begin{itemize}[nosep, leftmargin=1.2em]
    \item \textbf{Desviación total en horas}: $-$5\,h sobre 497\,h
    planificadas ($-$1\%). Dentro de tolerancia.
    \item \textbf{CPI estimado}: $\approx$1{,}01 (proyecto en
    presupuesto).
    \item \textbf{Contingencia consumida}: 0\% (no se necesitaron
    horas extra ponderadas por riesgo). El margen del 15\% queda como
    holgura.
    \item \textbf{Tareas opcionales canceladas}: T27 (benchmarks de
    latencia), redirigida a trabajo futuro.
    \item \textbf{Tareas no previstas incorporadas}: T20bis (logit
    lens), por necesidad metodológica detectada en G6.
\end{itemize}
